%% kapa: exam=1 round=30 year=2019 period=1 form=A subject=real_estate subject_ko=부동산학원론 seq=38 number=118
%% source: Q-Net 감정평가사 기출문제 (cst00309) artlSeq=5205205 file=extracted/1008/5205205/2019년 제30회 감정평가사 1차시험 문제지/2019년 제30회 감정평가사 1차시험 문제지/2019년도 제30회 감정평가사 1차 1교시 A형.hwp
%% conversion: kordoc 4.24.51 -> tools/gen1.py (auto); review=proofread-sonnet; flags=none
\begin{question}{118}
컨버스(P. Converse)의 분기점모형에 따르면 상권은 거리의 제곱에 반비례하고 인구에 비례한다. 다음의 조건에서 A, B 도시의 상권 경계지점은 A시로부터 얼마나 떨어진 곳에 형성되는가? (단, 주어진 조건에 한함)
\begin{framedtext}
○ A시의 인구: 16만명, B시의 인구: 4만명 \par
○ 두 도시 간의 거리: 15km \par
○ 두 도시의 인구는 모두 구매자이며, 두 도시에서만 구매함
\end{framedtext}
\begin{choices}
  \item 8km
  \item 9km
  \item 10km
  \item 11km
  \item 12km
\end{choices}
\end{question}
\answer{3}
